% ============================================================
% CHƯƠNG 4: HIỆN THỰC HỆ THỐNG
% MỤC ĐÍCH: Bổ sung chi tiết cho Câu 4 — chứng minh đã implement thực sự
% PHÂN BỔ: ~8–10 trang
% ============================================================
\chapter{Hiện thực hệ thống}
\label{chap:hien-thuc}

Trong khi Chương~\ref{chap:kien-truc} mô tả \emph{thiết kế ở mức contract}
giữa bốn lớp, chương này đi xuống mức \emph{mã nguồn thực thi} và trình
bày các quyết định hiện thực nhỏ hơn nhưng quyết định sự đúng đắn của hệ
thống: chiến lược parsing CSV chịu lỗi, kỹ thuật theo dõi tài nguyên thủ
công khi đồng hồ CloudSim chưa khởi động, đồng bộ hai luồng qua
\texttt{SynchronousQueue}, chuyển đổi mảng Java\,$\to$\,\texttt{numpy} an
toàn, và xử lý các sự cố thực tế trong quá trình tích hợp. Toàn bộ Phase 1
hiện có $3\,188$ dòng mã chia giữa $10$ lớp Java
($\approx 1\,800$ dòng) và $7$ module Python ($\approx 1\,400$ dòng);
Bảng~\ref{tab:loc-summary} liệt kê cụ thể.

\begin{table}[H]
  \centering
  \caption{Quy mô mã nguồn Phase 1 (đếm bằng \texttt{wc -l}, không trừ
  comment/blank line). Tổng $3\,188$ dòng.}
  \label{tab:loc-summary}
  \footnotesize
  \begin{tabularx}{\linewidth}{l X r}
    \toprule
    \textbf{Module} & \textbf{Mô tả ngắn} & \textbf{LOC} \\
    \midrule
    \texttt{SimulationConfig.java}        & Cấu hình tập trung host/power/QoS                & 80   \\
    \texttt{DatacenterFactory.java}       & Tạo Datacenter + GPU registry                    & 153  \\
    \texttt{AlibabaTraceReader.java}      & Parser CSV chịu lỗi $\to$ \texttt{TaskRecord}    & 160  \\
    \texttt{ScenarioFilter.java}          & Lọc Low/High/Burst                               & 122  \\
    \texttt{MetricsExporter.java}         & Ghi CSV + JSON                                   & 164  \\
    \texttt{SimulationManager.java}       & Vòng đời + RL stepping (CRITICAL)                & 492  \\
    \texttt{VmAllocPolicyK8sDefault.java} & Filter + LeastRequestedPriority                  & 124  \\
    \texttt{VmAllocPolicyRandom.java}     & Random uniform trên tập khả thi                  & 81   \\
    \texttt{GatewayEntryPoint.java}       & API Py4J + bootstrap GatewayServer               & 247  \\
    \texttt{Main.java}                    & Entrypoint daemon                                & 31   \\
    \midrule
    \texttt{state\_builder.py}            & \texttt{double[]}\,$\to$\,\texttt{ndarray} + parse & 129  \\
    \texttt{reward.py}                    & Welford normalizer + scalarise                   & 172  \\
    \texttt{environment.py}               & \texttt{gym.Env} + \texttt{ScalarRewardWrapper}  & 374  \\
    \texttt{tracker.py}                   & WandB online/offline                             & 243  \\
    \texttt{baseline\_eval.py}            & Đánh giá K8s/Random qua Py4J                     & 242  \\
    \texttt{smoke\_test.py}               & End-to-end test 8 phase                          & 350  \\
    \texttt{train.py}                     & Stub Phase 2                                     & 24   \\
    \bottomrule
  \end{tabularx}
\end{table}

% ============================================================
\section{Hiện thực phân hệ Java}
\label{sec:impl-java}

\subsection{Cấu hình dự án Maven}
\label{ssec:maven}

Dự án Java dùng \emph{Maven 3.9.9} với hai dependency chính: CloudSim
Plus 8.5.7 và Py4J 0.10.9.7. Listing~\ref{lst:pom} trích đoạn quan trọng
của \texttt{pom.xml}; điểm đáng chú ý là:
\begin{itemize}
  \item \textbf{ArtifactId chính xác.} Phải dùng \texttt{cloudsimplus}
        (không có dấu gạch ngang) -- artifact \texttt{cloudsim-plus} cũ
        đã bị deprecate kể từ phiên bản 8 và không còn trên Maven Central.
        Sai sót này từng tốn $\sim 2$ giờ debug khi Maven báo
        \texttt{Could not resolve dependencies}.
  \item \textbf{\texttt{maven-shade-plugin}} đóng gói toàn bộ phụ thuộc
        thành một \emph{fat JAR} (\texttt{simulation-1.0-SNAPSHOT.jar},
        $\approx 18$\,MB), giúp container runtime chỉ cần một lệnh
        \texttt{java -jar} mà không cần \texttt{CLASSPATH} ngoài.
  \item \textbf{\texttt{<release>21</release>}} đặt tường minh cho
        \texttt{maven-compiler-plugin}; CloudSim Plus 8.5.7 sử dụng
        Java \texttt{record}, \texttt{sealed} và \texttt{pattern matching
        for switch} -- yêu cầu tối thiểu Java 21.
\end{itemize}

\begin{lstlisting}[style=javastyle, caption={\texttt{pom.xml} -- trích đoạn
hai dependency chính.}, label={lst:pom},
language={[LaTeX]TeX}]
<dependency>
  <groupId>org.cloudsimplus</groupId>
  <artifactId>cloudsimplus</artifactId>   <!-- NOT "cloudsim-plus" -->
  <version>8.5.7</version>
</dependency>
<dependency>
  <groupId>net.sf.py4j</groupId>
  <artifactId>py4j</artifactId>
  <version>0.10.9.7</version>             <!-- pinned: matches Python -->
</dependency>
\end{lstlisting}

\subsection{\texttt{DatacenterFactory} -- xây dựng datacenter và registry GPU}
\label{ssec:impl-factory}

Theo thiết kế ở §\ref{ssec:factory}, mỗi \texttt{Host} CloudSim được kèm
một \texttt{GpuState} song song trong một \texttt{IdentityHashMap}. Hàm
factory \texttt{create(simulation, dcSpec, gpuRegistry)} thực hiện sáu
bước:

\begin{enumerate}
  \item Sinh \texttt{HostSpec.pesCount()} đối tượng \texttt{PeSimple} cho
        từng host;
  \item Khởi tạo \texttt{HostSimple} với
        \texttt{ResourceProvisionerSimple} cho RAM/băng thông;
  \item Cấu hình \texttt{VmSchedulerTimeShared} -- chia thời gian PE giữa
        các Vm;
  \item Gắn \texttt{PowerModelHostSimple(maxPower, idlePower)} áp công
        thức~\eqref{eq:linear-power};
  \item Tạo \texttt{GpuState(gpuCount)} mới và đặt vào \texttt{gpuRegistry}
        với chính \texttt{Host} làm khoá (so sánh tham chiếu);
  \item Trả về \texttt{DatacenterSimple} bao gói toàn bộ.
\end{enumerate}

Hàm \texttt{gpuPowerWatt(gpuState, powerSpec)} hiện thực
công thức~\eqref{eq:gpu-power}: tổng công suất GPU của host bằng
$n\cdot P^{\mathrm{idle}} + g\cdot (P^{\mathrm{max}} - P^{\mathrm{idle}})$
với $g$ là số GPU đang dùng, được \texttt{SimulationManager} gọi mỗi bước
để tính reward $r^{\mathrm{energy}}_t$.

\subsection{\texttt{AlibabaTraceReader} -- parser chịu lỗi}
\label{ssec:impl-trace}

File trace có $\sim 8\,000$ dòng; trace thực tế \emph{không} sạch tuyệt
đối: một số dòng có trường rỗng, giá trị \texttt{nan}, hoặc số cột thiếu
do ký tự \texttt{,} bên trong tên tác vụ. Triết lý của parser là
\textbf{tolerant by design}: bỏ qua dòng lỗi và in cảnh báo, không huỷ
toàn bộ import vì $20$\,dòng hỏng trong $8\,000$ là chấp nhận được.

\begin{lstlisting}[style=javastyle, caption={\texttt{AlibabaTraceReader.read()}
-- tolerant CSV parsing.}, label={lst:trace-read}]
public static List<TaskRecord> read(String csvPath) throws IOException {
    Path path = Path.of(csvPath);
    List<TaskRecord> records = new ArrayList<>();
    int skipped = 0;
    try (BufferedReader br = Files.newBufferedReader(path)) {
        br.readLine();                       // skip header
        String line;  int lineNum = 1;
        while ((line = br.readLine()) != null) {
            lineNum++;
            TaskRecord rec = parseLine(line, lineNum);
            if (rec != null) records.add(rec);
            else             skipped++;
        }
    }
    records.sort(Comparator.comparingDouble(TaskRecord::creationTime));
    System.out.printf("[AlibabaTraceReader] Loaded %d tasks from %s "
                    + "(skipped %d)%n", records.size(), csvPath, skipped);
    return Collections.unmodifiableList(records);
}
\end{lstlisting}

Hai helper \texttt{parseIntSafe} và \texttt{parseDoubleSafe} xử lý chuỗi
rỗng và \texttt{nan} bằng giá trị fallback ($0$); ngoại lệ
\texttt{NumberFormatException} bị bắt ở mức từng dòng để parser không
phải xử lý cấp thấp hơn. \emph{Deadline} của tác vụ được dẫn xuất bằng
chính \texttt{deletionTime} -- giá trị này đại diện cho thời điểm tác vụ
\emph{thực sự kết thúc} trong K8s gốc, do đó là tham chiếu hợp lý cho
một scheduler thay thế. Hệ số $\lambda$ được nạp ngay tại đây qua
\texttt{SimulationConfig.qosToLambda(qos)}.

\subsection{Hai baseline schedulers}
\label{ssec:impl-baselines}

Cả hai baseline kế thừa \texttt{VmAllocationPolicyAbstract} của CloudSim
Plus và \emph{đồng thời} cung cấp một entry-point độc lập
\texttt{selectHostForTask(hosts, task)} dành cho con đường evaluation
qua Py4J. Lý do tách hai con đường: pipeline mô phỏng nội bộ của CloudSim
chỉ làm việc với \texttt{Vm} (không có thông tin GPU) trong khi vòng RL
phải biết đầy đủ về tác vụ, kể cả số GPU yêu cầu. Hai con đường này
phối hợp ăn khớp nhau qua cùng tập \texttt{gpuRegistry} dùng chung.

\textbf{Kubernetes-default scheduler.} Listing~\ref{lst:k8s-impl} là
phần lõi của \texttt{VmAllocationPolicyK8sDefault}: lọc, chấm
\textit{LeastRequestedPriority} (LRP), và tie-break bằng chỉ số nhỏ nhất
để đảm bảo tất định.

\begin{lstlisting}[style=javastyle, caption={\texttt{VmAllocationPolicyK8sDefault}
-- Filter $\to$ Score $\to$ Select.}, label={lst:k8s-impl}]
public int selectHostForTask(List<Host> hosts, TaskRecord task) {
    int    bestIdx = -1;
    double bestScore = -1;
    for (int i = 0; i < hosts.size(); i++) {
        Host host = hosts.get(i);
        if (!canHost(host, task)) continue;            // Filter
        double score = leastRequestedScore(host);      // Score
        if (score > bestScore) {                       // Select (deterministic)
            bestScore = score;
            bestIdx   = i;
        }
    }
    return (bestIdx >= 0) ? bestIdx : 0;
}

private double leastRequestedScore(Host host) {
    double cpuFrac = (double) host.getFreePesNumber() / host.getPesNumber();
    double memFrac = (double) host.getRam().getAvailableResource()
                              / host.getRam().getCapacity();
    return (cpuFrac + memFrac) / 2.0;                  // K8s LRP
}
\end{lstlisting}

\textbf{Random scheduler.} Sàng lọc tập khả thi rồi lấy ngẫu nhiên
đồng đều với \texttt{Random(seed)} -- seed cùng \texttt{RANDOM\_SEED}
chung của hệ thống đảm bảo \emph{mỗi lần chạy lại trên cùng kịch bản đều
cho cùng một dãy hành động ngẫu nhiên}. Đây là yêu cầu cứng để bảng số
liệu Phase 2 có thể được tái sản xuất.

\subsection{\texttt{SimulationManager} -- chi tiết kỹ thuật quan trọng}
\label{ssec:impl-sim-manager}

\paragraph{Tracking thủ công thay vì \texttt{simulation.start()}.}
Như đã đề cập ở §\ref{ssec:sim-manager}, Phase 1 \emph{cố ý} không khởi
động đồng hồ CloudSim. Hệ quả: \texttt{Host.getCpuPercentUtilization()}
luôn trả $0$, \texttt{getRam().getAvailableResource()} luôn trả full.
Để có dữ liệu thật cho observation và reward, \texttt{SimulationManager}
duy trì hai bản đồ song song:

\begin{lstlisting}[style=javastyle, caption={Manual resource accounting.
Đây là \emph{nguồn sự thật} duy nhất cho utilisation trong Phase 1.},
label={lst:manual-tracking}]
private final Map<Host, Integer> hostPeUsage  = new HashMap<>();
private final Map<Host, Long>    hostRamUsage = new HashMap<>();

private void allocateTask(TaskRecord task, Host host) {
    hostPeUsage.merge(host, task.pesNeeded(),       Integer::sum);
    hostRamUsage.merge(host, (long) task.memoryMib(), Long::sum);
    if (task.numGpu() > 0) {
        GpuState g = gpuRegistry.get(host);
        if (g != null) g.allocate(task.numGpu());
    }
}
\end{lstlisting}

Quyết định này được \emph{lock-in} qua một dòng comment chỉ dẫn rõ
ràng trong source và được nhắc lại trong \texttt{CLAUDE.md} để các phiên
phát triển sau không vô tình ``sửa cho đúng'' (gọi
\texttt{simulation.start()}) -- vốn sẽ phát sinh sự kiện không có handler
và crash JVM âm thầm.

\paragraph{Đồng bộ hai luồng qua \texttt{SynchronousQueue}.}
\texttt{stepResultQueue} và \texttt{actionQueue} là hai
\texttt{SynchronousQueue<T>}: \emph{không có buffer}. Một lời gọi
\texttt{put(x)} chỉ trả về khi một bên \texttt{take()} đang sẵn sàng nhận.
Điều này biến mỗi cặp $(\texttt{step()},\,\texttt{stepResultQueue.put})$
thành một giao dịch nguyên tử nhịp đôi với phía Python -- không có race
condition giữa hai luồng mà không cần \texttt{synchronized}. Khi reset
giữa hai episode, hàm \texttt{terminateExisting()} ngắt thread cũ bằng
\texttt{interrupt() + join(2000ms)} rồi \texttt{poll()} drain hai queue
để loại bỏ dữ liệu còn sót -- một mẫu thiết kế chuẩn cho thread
producer/consumer trong Java.

\paragraph{Tính reward đồng nhất với phát biểu hình thức.}
Hàm \texttt{computeReward} (Listing~\ref{lst:reward-java} ở Chương~3)
hiện thực đúng cặp công thức~\eqref{eq:reward-step}: $r^{\mathrm{energy}}_t
= -\Delta E_t$ và $r^{\mathrm{SLA}}_t = -\lambda(q_i)\,\max(0,\,\hat C_i
- d_i)$. Khi $\hat C_i > d_i$, đồng thời tăng \texttt{slaViolationCount}
để \texttt{MetricsExporter} sau này tính tỉ lệ vi phạm.

\subsection{\texttt{MetricsExporter} -- xuất CSV + JSON không phụ thuộc}
\label{ssec:impl-export}

\texttt{MetricsExporter} chủ ý \emph{không thêm dependency JSON}
(như Jackson hay Gson) -- Java 15+ đã có \emph{text block}
(\texttt{"""\,...\,"""}) đủ để format JSON ngắn gọn và rõ ràng. Hai
\texttt{record} tài liệu hoá schema:

\begin{itemize}
  \item \texttt{Snapshot}: $9$ trường tương ứng $9$ cột trong
        \texttt{metrics.csv} -- một dòng mỗi tác vụ. Pandas có thể đọc
        thẳng bằng \texttt{pd.read\_csv}.
  \item \texttt{Summary}: $9$ trường tổng kết episode -- ghi ra
        \texttt{summary.json} với khoá \texttt{snake\_case} thân thiện
        Python.
\end{itemize}

Schema CSV:
\begin{verbatim}
timestamp, cpu_energy_kwh, gpu_energy_kwh, total_energy_kwh,
tasks_scheduled, sla_violations, avg_sla_slack,
cluster_cpu_util, cluster_gpu_util
\end{verbatim}
Lựa chọn \texttt{Locale.US} cho \texttt{String.format} đảm bảo dấu thập
phân là \texttt{.} bất kể locale của máy chủ -- tránh trường hợp
\texttt{,} của locale Việt Nam làm hỏng cấu trúc CSV.

% ============================================================
\section{Hiện thực phân hệ Python}
\label{sec:impl-python}

\subsection{\texttt{state\_builder.py} -- chuyển đổi an toàn và \emph{parsable}}
\label{ssec:impl-state}

Hàm \texttt{from\_java(java\_array, num\_hosts)} làm ba việc theo thứ tự:
(i)~chuyển \texttt{double[]} (qua Py4J) thành
\texttt{numpy.float32}; (ii)~validate độ dài đúng bằng $3H + 4$, raise
\texttt{ValueError} có thông điệp tường minh nếu không khớp;
(iii)~\texttt{np.clip(obs, 0.0, 1.0, out=obs)} -- chống lỗi tràn nhỏ do
phía Java đôi khi trả $> 1{,}0$ trong chế độ time-shared. Sau đó vector
này được trả thẳng ra Gymnasium mà không phải copy thêm.

Để tiện debug, module còn cung cấp \texttt{dataclass(frozen=True,
slots=True)} \texttt{ParsedObservation} cho phép truy cập các thành phần
theo tên thay vì chỉ số. Dùng \texttt{slots=True} để giảm chi phí bộ nhớ
khi log hàng triệu observation trong huấn luyện Phase 2.

\subsection{\texttt{reward.py} -- vector, normalisation, scalarisation}
\label{ssec:impl-reward}

Module này có ba thành phần độc lập, có thể bật/tắt riêng:
\begin{enumerate}
  \item \texttt{from\_java(jrew)}: kiểm tra độ dài $= 2$, ép kiểu
        \texttt{float32}, không clip vì reward âm có ý nghĩa.
  \item \texttt{RewardNormalizer}: chuẩn hoá Welford online từng thành
        phần độc lập, công thức cập nhật trung bình $\mu$ và phương sai
        $M_2$:
        \[
          \mu_n = \mu_{n-1} + \frac{x_n - \mu_{n-1}}{n}, \qquad
          M_{2,n} = M_{2,n-1} + (x_n - \mu_{n-1})(x_n - \mu_n).
        \]
        Khi cần dùng (\texttt{normalize\_reward=True}), wrapper trả
        $(x - \mu)/\sigma$ với $\sigma = \sqrt{M_2 / (n-1) + \epsilon}$,
        $\epsilon = 10^{-8}$. Mục đích: ngăn $r^{\mathrm{energy}}$
        ($\sim -300$) lấn át $r^{\mathrm{SLA}}$ ($\sim 0$ phần lớn bước)
        trong giai đoạn đầu PPO.
  \item \texttt{scalarise(rvec, weights)}: tích vô hướng
        $\langle \mathbf{w},\,\mathbf{r}\rangle$, hiện thực~\eqref{eq:linear-scalar}.
        \texttt{weights} mặc định nạp từ \texttt{ENERGY\_WEIGHT} và
        \texttt{SLA\_WEIGHT} trong \texttt{.env}, được chuẩn hoá cho
        $w_{\mathrm{e}} + w_{\mathrm{s}} = 1$.
\end{enumerate}

\subsection{\texttt{environment.py} -- gắn kết với Gymnasium API}
\label{ssec:impl-env}

\texttt{CloudSimEnv} kế thừa \texttt{gym.Env} và tuân thủ chặt giao diện
Gymnasium 1.0+: \texttt{reset()} trả \texttt{(obs, info)};
\texttt{step()} trả \texttt{(obs, reward, terminated, truncated, info)}.
Có hai chi tiết hiện thực đáng lưu ý:

\paragraph{Reset bootstrap.} Khi khởi tạo, \texttt{\_\_init\_\_}
\emph{phải gọi reset một lần} mới biết được số host (do
\texttt{getActionSize()} chỉ hợp lệ sau \texttt{resetSimulation()}).
Để tránh tốn một episode reset thừa, kết quả này được cache vào
\texttt{self.\_initial\_result} và tái sử dụng cho lần \texttt{reset()}
đầu tiên nếu seed/scenario không đổi:

\begin{lstlisting}[style=pythonstyle, caption={\texttt{CloudSimEnv.reset}
-- tránh reset thừa lần đầu.}, label={lst:env-reset}]
def reset(self, *, seed=None, options=None):
    ep_seed     = seed if seed is not None else self._seed
    ep_scenario = (options or {}).get("scenario", self._scenario)
    if (self._needs_initial_reset
            and ep_seed == self._seed
            and ep_scenario == self._scenario):
        result = self._initial_result          # reuse __init__ reset
        self._needs_initial_reset = False
    else:
        result = self._ep.reset(ep_scenario, ep_seed)
    self._current_obs = state_builder.from_java(result.observation(),
                                                 self._num_hosts)
    self._done = False
    self._step_count = 0
    self._episode_rewards = []
    return self._current_obs, {"task_index": int(result.taskIndex()),
                               "task_name":  str(result.taskName()),
                               "scenario":   ep_scenario,
                               "seed":       ep_seed}
\end{lstlisting}

\paragraph{Action mask an toàn.} Hàm \texttt{action\_masks()} (gọi
sang Java) có thể trả mask toàn \texttt{False} nếu không host nào còn
chứa nổi tác vụ. Trong trường hợp đó, Python bật toàn bộ mask về
\texttt{True} để episode không bị khoá -- Java sau đó vẫn có quyền
phạt nặng qua $r^{\mathrm{SLA}}$, nên bài toán không bị méo:

\begin{lstlisting}[style=pythonstyle, caption={\texttt{action\_masks()}
-- fail-open để không khoá episode.}, label={lst:action-masks}]
def action_masks(self) -> np.ndarray:
    java_mask = self._ep.getActionMask()
    mask = np.array(list(java_mask), dtype=bool)
    if not mask.any():            # safety: avoid stuck episode
        mask[:] = True
    return mask
\end{lstlisting}

\paragraph{\texttt{ScalarRewardWrapper}.} Phụ thuộc trọng số quét trong
Phase 2, lớp wrapper kế thừa \texttt{gym.RewardWrapper} chỉ cần override
\texttt{reward(self, reward)} và truyền \texttt{action\_masks} ngược về
môi trường gốc -- giữ nguyên hợp đồng MaskablePPO.

\paragraph{Đóng kết nối sạch.} \texttt{close()} gọi
\texttt{gateway.shutdown(send\_shutdown=False)} để đóng kết nối
\emph{phía client} mà không tắt \texttt{GatewayServer} của Java -- điều
này cho phép \texttt{baseline\_eval.py} chạy lần lượt nhiều scheduler
trên cùng một tiến trình Java mà không phải khởi động lại JVM (chi phí
$\sim 3$\,giây mỗi lần).

\subsection{\texttt{tracker.py} -- WandB online/offline tự động}
\label{ssec:impl-tracker}

Trong điều kiện phòng thí nghiệm không phải lúc nào cũng có Internet
hoặc \texttt{WANDB\_API\_KEY}. \texttt{ExperimentTracker} dùng pattern
\emph{degrade gracefully}: hàm \texttt{\_should\_use\_wandb()} kiểm tra
hai điều kiện -- thư viện \texttt{wandb} có cài và \texttt{WANDB\_API\_KEY}
không rỗng -- nếu cả hai đúng thì online, ngược lại fallback console.
Mọi phương thức \texttt{log\_step}, \texttt{log\_episode},
\texttt{log\_comparison}, \texttt{save\_model} đều có nhánh fallback
console riêng, cùng schema khoá \texttt{step/...} và \texttt{episode/...}
để dashboard WandB có thể bật mà không cần đổi code agent.

\subsection{\texttt{baseline\_eval.py} -- so sánh fair}
\label{ssec:impl-baseline-eval}

Để đảm bảo tính \emph{fair} khi so sánh ba scheduler ở Phase 2, baseline
phải vận hành \emph{trong cùng môi trường mô phỏng} với agent RL.
\texttt{baseline\_eval.py} đạt được điều này bằng cách \emph{không}
hiện thực lại logic K8s/Random ở Python; thay vào đó nó gọi
\texttt{env.\_ep.selectBaselineAction("k8s")} để hỏi Java chính bộ
\texttt{VmAllocationPolicy} đã lock seed (Listing~\ref{lst:baseline-loop}).
Nhờ vậy, mọi sự khác biệt giữa Random và MORL chỉ phản ánh chất lượng
chính sách, không phải chi tiết hiện thực.

\begin{lstlisting}[style=pythonstyle, caption={\texttt{baseline\_eval.evaluate\_baseline}
-- tận dụng chung pipeline với RL.}, label={lst:baseline-loop}]
def evaluate_baseline(env, policy_name, scenario="HIGH", seed=42):
    obs, info = env.reset(seed=seed, options={"scenario": scenario})
    total_e, total_s, steps = 0.0, 0.0, 0
    while True:
        action = int(env._ep.selectBaselineAction(policy_name))   # ask Java
        obs, rvec, terminated, truncated, info = env.step(action)
        total_e += float(rvec[0]);  total_s += float(rvec[1]);  steps += 1
        if terminated or truncated:
            break
    return {"steps": steps, "total_energy_reward": total_e,
            "total_sla_reward": total_s,
            "total_energy_kwh": float(env._ep.getTotalEnergyKwh())}
\end{lstlisting}

% ============================================================
\section{Hiện thực hạ tầng Docker}
\label{sec:impl-docker}

\subsection{Pipeline build hai stage}
\label{ssec:impl-build}

Cả hai \texttt{Dockerfile} đều dùng pattern multi-stage; số liệu kích
thước image đo được trên máy phát triển (Windows 11, Docker Desktop
4.x) liệt kê ở Bảng~\ref{tab:image-size}.

\begin{figure}[H]
  \centering
  \includegraphics[width=0.9\textwidth]{screenshot-build-success.png}
  \caption{Kết quả \texttt{docker compose build} hoàn tất cho cả hai
  service trên máy phát triển.}
  \label{fig:build-success}
\end{figure}

\begin{table}[H]
  \centering
  \caption{Kích thước Docker image cuối cùng}
  \label{tab:image-size}
  \footnotesize
  \begin{tabular}{lccr}
    \toprule
    \textbf{Service} & \textbf{Base runtime} & \textbf{Stage builder} & \textbf{Cuối cùng} \\
    \midrule
    \texttt{eldas-cloudsim-java} & \texttt{eclipse-temurin:21-jre-jammy} & 1\,1xx\,MB (JDK+Maven) & \textbf{$\sim 434$\,MB} \\
    \texttt{eldas-rl-agent}      & \texttt{python:3.11.9-slim-bookworm}   & $\sim 700$\,MB (gcc)   & \textbf{$\sim 8{,}57$\,GB} \\
    \bottomrule
  \end{tabular}
\end{table}

\textbf{Cache layer.} Trick \emph{copy manifest trước source} đem lại
hiệu quả thực tế đo được: rebuild đầy đủ ($1$--$2$\,phút) so với rebuild
sau khi chỉ sửa source ($\sim 10$\,giây). Cụ thể trong Dockerfile Java:

\begin{lstlisting}[style=yamlstyle, caption={Trick layer cache: \texttt{pom.xml}
trước, \texttt{src/} sau.}, label={lst:dockerfile-cache}]
FROM eclipse-temurin:21-jdk-jammy AS builder
ARG MAVEN_VERSION=3.9.9
RUN apt-get update && apt-get install -y --no-install-recommends wget \
    && wget -q https://archive.apache.org/dist/maven/maven-3/${MAVEN_VERSION}/...
WORKDIR /build
COPY pom.xml .                                  # layer 1: stable
RUN mvn dependency:go-offline -B --no-transfer-progress
COPY src ./src                                  # layer 2: changes often
RUN mvn package -DskipTests -B --no-transfer-progress

FROM eclipse-temurin:21-jre-jammy AS runtime    # ~190MB JRE
COPY --from=builder /build/target/simulation-1.0-SNAPSHOT.jar /app/sim.jar
\end{lstlisting}

\subsection{Quản lý volume}
\label{ssec:impl-volumes}

Ba volume được khai báo trong \texttt{docker-compose.yml} với chiến
lược khác nhau theo vai trò:
\begin{itemize}
  \item \texttt{trace-data}: \emph{bind mount} từ
        \texttt{./data/alibaba-trace} của host. Mount \texttt{:ro}
        cho \texttt{cloudsim-java} -- Java không thể vô tình ghi đè.
  \item \texttt{results}: bind mount từ \texttt{./data/results}, được
        \emph{cả hai service} gắn (Java ghi
        \texttt{metrics.csv}/\texttt{summary.json}, Python đọc lại để
        log lên WandB).
  \item \texttt{model-store}: \emph{named volume} do Docker quản lý,
        chỉ \texttt{rl-agent} truy cập, lưu checkpoint mô hình PPO.
\end{itemize}
Phân tách rõ \emph{data đầu vào} (read-only) vs \emph{output} (read-write
chia sẻ) vs \emph{artifact riêng} (Docker-managed) là ứng dụng nguyên tắc
\textit{principle of least privilege} ở mức storage.

\subsection{Vấn đề kỹ thuật gặp phải và giải pháp}
\label{ssec:impl-issues}

Bảng~\ref{tab:issues} tổng hợp \emph{sáu sự cố thực tế} đã gặp trong
quá trình hiện thực Phase 1, kèm nguyên nhân gốc rễ và giải pháp đã áp
dụng. Đây không phải checklist lý thuyết -- mỗi mục đều tương ứng với
một commit cụ thể trong lịch sử Git.

\begin{table}[H]
  \centering
  \caption{Sáu vấn đề kỹ thuật đáng ghi nhận trong Phase 1 và giải pháp.}
  \label{tab:issues}
  \footnotesize
  \begin{tabularx}{\linewidth}{l X X}
    \toprule
    \textbf{Vấn đề} & \textbf{Nguyên nhân} & \textbf{Giải pháp đã áp dụng} \\
    \midrule
    Maven không tải được \texttt{cloudsim-plus}
      & Artifact ID đã đổi tên sang \texttt{cloudsimplus} (không gạch ngang) từ v8.
      & Sửa \texttt{pom.xml} dùng \texttt{cloudsimplus}; tài liệu hoá trong
        \texttt{CLAUDE.md} để không lặp lại. \\
    \texttt{Host.getCpuPercentUtilization()} luôn bằng $0$
      & \texttt{simulation.start()} chưa được gọi $\Rightarrow$
        đồng hồ CloudSim chưa chạy.
      & Bypass: theo dõi tài nguyên thủ công bằng
        \texttt{HashMap<Host, Integer>}; lock-in bằng comment + CLAUDE.md. \\
    JVM crash âm thầm khi gọi \texttt{broker.submitVm}
      & \texttt{simulation.schedule()} đẩy event vào loop chưa khởi động
        $\Rightarrow$ NPE deep inside CloudSim.
      & Không gọi \texttt{submitVm} ở Phase 1; allocate qua bookkeeping
        thủ công (Listing~\ref{lst:manual-tracking}). \\
    \texttt{rl-agent} kết nối Py4J fail ngay sau \texttt{up}
      & Healthcheck mark Java healthy ngay khi port mở, nhưng
        \texttt{GatewayServer.bind} chậm thêm $\sim 2$\,giây.
      & Lớp phòng vệ thứ hai: retry $10\times$ với backoff $3$\,s ở
        \texttt{environment.\_connect\_gateway}. \\
    Xung đột phiên bản \texttt{gymnasium}
      & \texttt{stable-baselines3 2.7.0} cần
        \texttt{gymnasium\,>=\,1.0.0}; \texttt{mo-gymnasium 1.3.2} cần
        \texttt{gymnasium\,>=\,0.28}.
      & \emph{Không pin} \texttt{gymnasium} trong
        \texttt{requirements.txt}; để pip tự giải xung đột. \\
    Race condition giữa simThread và Python
      & Hai luồng chia sẻ \texttt{currentTaskIdx} không có lock.
      & Dùng \texttt{SynchronousQueue} (no buffer) -- biến mỗi step
        thành giao dịch nguyên tử mà không cần \texttt{synchronized}. \\
    \bottomrule
  \end{tabularx}
\end{table}

% ============================================================
\section{Smoke test end-to-end}
\label{sec:impl-smoke}

Để xác nhận bốn lớp nối đúng nhau, ELDAS kèm \texttt{smoke\_test.py}
($350$ dòng) chạy tám pha kiểm thử tăng dần độ phủ
(Bảng~\ref{tab:smoke-phases}). Mỗi pha là một \texttt{assert} cô lập
với thông báo lỗi rõ ràng -- nếu pha $k$ thất bại, các pha $> k$ không
chạy.

\begin{table}[H]
  \centering
  \caption{Tám pha của smoke test \texttt{T4.5}.}
  \label{tab:smoke-phases}
  \footnotesize
  \begin{tabularx}{\linewidth}{l X}
    \toprule
    \textbf{Pha} & \textbf{Kiểm tra} \\
    \midrule
    1 -- Connection      & Py4J kết nối được, \texttt{getActionSize > 0},
                           \texttt{observation\_space} đúng $(3H+4,)$ \\
    2 -- Reset           & \texttt{reset()} trả \texttt{obs} đúng shape,
                           dtype \texttt{float32}, mọi giá trị $\in [0,1]$ \\
    3 -- Action mask     & Mask đúng shape $(H,)$, dtype \texttt{bool},
                           có ít nhất một host khả thi \\
    4 -- Single step     & \texttt{step()} trả 5-tuple đúng kiểu;
                           $r^{\mathrm{energy}} \le 0$, $r^{\mathrm{SLA}} \le 0$ \\
    5 -- Full episode    & Chạy hết tác vụ; tự nhiên \texttt{terminated=True};
                           kiểm tra $r^{\mathrm{energy}} < -120$
                           (ngưỡng idle power) cho chứng cứ ``manual PE
                           tracking đang hoạt động'' \\
    6 -- Metrics export  & \texttt{exportMetrics(\$DIR)} ghi
                           \texttt{metrics.csv} + \texttt{summary.json} \\
    7 -- Scalar wrapper  & \texttt{ScalarRewardWrapper} trả \texttt{float};
                           giá trị $\approx \mathbf{w}\cdot\mathbf{r}$ \\
    8 -- Baseline action & \texttt{selectBaselineAction("k8s")} và
                           \texttt{("random")} trả chỉ số host hợp lệ \\
    \bottomrule
  \end{tabularx}
\end{table}

\paragraph{Cách chạy.}
Sau khi \texttt{docker compose up --build} thấy cả hai service \texttt{healthy}:
\begin{lstlisting}[style=yamlstyle, caption={Lệnh chạy smoke test.}, label={lst:run-smoke}]
docker compose run --rm rl-agent python src/smoke_test.py \
    --scenario LOW --max-steps 3000 --output /data/results/smoke
\end{lstlisting}

Kết quả mong đợi: tám pha đều \texttt{[OK]}, tổng thời gian
$\sim 5$--$8$\,giây trên CPU thường, đầu ra
\texttt{data/results/smoke/metrics.csv} (\texttt{$\sim 2\,000$ dòng}) và
\texttt{summary.json}. Chương~\ref{chap:ket-qua} sẽ trình bày kết quả
cụ thể: bảng giá trị reward/energy, ảnh chụp log, ảnh chụp file CSV.

% ============================================================
\section{Tổng kết chương}
\label{sec:ket-luan-c4}

Chương này đã đi qua từng quyết định hiện thực ở mức mã nguồn của bốn
lớp:
\begin{itemize}
  \item \textbf{Lớp 2 (Java).} Maven + fat JAR, parser CSV chịu lỗi,
        hai baseline scheduler chia sẻ \texttt{gpuRegistry}, manual
        resource tracking, đồng bộ hai luồng qua \texttt{SynchronousQueue},
        và xuất metrics không phụ thuộc.
  \item \textbf{Lớp 4 (Python).} Chuyển đổi an toàn Java\,$\to$\,numpy,
        Welford normaliser, wrapper Gymnasium chuẩn 1.0+, fallback
        WandB tự động.
  \item \textbf{Lớp 1 (Docker).} Multi-stage build, layer cache trick,
        bind mount read-only cho dữ liệu đầu vào, retry phía client
        bù healthcheck timing.
  \item \textbf{Smoke test} kiểm chứng cả vòng RL trên cùng pipeline
        đã sẵn sàng cho huấn luyện Phase 2.
\end{itemize}

Bảng~\ref{tab:issues} đã tổng hợp \emph{sáu sự cố thực tế} -- bằng
chứng cho thấy hệ thống không phải lý thuyết suông mà đã đi qua đủ
chu kỳ \emph{thử-sai-sửa}. Chương kế tiếp trình bày kết quả định lượng
của Phase 1: log smoke test thành công, metrics
\texttt{energy/SLA} ban đầu của hai baseline trên ba kịch bản, và xác
nhận tính lặp lại của hệ thống qua \texttt{RANDOM\_SEED}.
